
\documentclass[a4paper,12pt]{article}

%Class
	%\documentclass[a4paper,12pt]{article}

%Packages
	\usepackage{amsmath}					% Standard maths
	\usepackage{amssymb}					% Standard maths
	\usepackage{amsthm}						% Standard maths
	\usepackage{thmtools, thm-restate}% Theorem styles
	\usepackage{mathtools}
	\usepackage[newzealand]{babel}% Date/time in British format (yes, I know it says NZ)
	\usepackage{datetime}					% Date/time
	\usepackage{multirow}					% For stacked subscripts
	\usepackage{xcolor}						% For colour!
	\usepackage[normalem]{ulem}		% For strikeout
	\usepackage{isomath}
	\usepackage{fancyhdr}					% For the headers
	\usepackage{mfirstuc}					% For the headers
	\usepackage[hmargin=1in,top=0.8in,bottom=1in]{geometry} % Page margins more reasonable
	\usepackage{graphicx}					% For including graphics
	%\usepackage{float}						% For something
	\usepackage{enumitem}					% For numbering tweaks
	\usepackage[margin=15pt,font={small},labelfont=bf,format=hang]{caption} % Changing caption style
	\usepackage{url}
	\usepackage[hidelinks]{hyperref}	% For hyperlinks
	\usepackage{breakurl}
	\usepackage{cleveref}
	\usepackage[scaled=.92]{helvet}
	\usepackage[]{FiraSans}
	\usepackage[T1]{fontenc}
	\usepackage{comment}
	\newbool{show_question}
	\newbool{show_solution}
	\newbool{show_commentary}
	\newcommand{\solutioninheading}{\ifbool{show_solution}{: SOLUTIONS\ifbool{show_commentary}{ \& FEEDBACK}{}}{}}
	
	\usepackage{pgfplots}
\newcommand{\ep}{\varepsilon}	

    \newcommand{\question}[3]{
    	\item %\textbf{#1} \par \medskip % Title
    	%
    	\ifbool{show_question}{#2}{} % Question
    	%
    	%
    	\ifbool{show_question}{\ifbool{show_solution}{\par\medskip\textbf{Solution:}}{}}{} % Both
    	\ifbool{show_solution}{#3}{} % Answer
    }

\usepackage[most]{tcolorbox}

%MATHEMATICAL SHORTCUTS
%Two and three-part definitions
	\newcommand{\twopartdef}[4]
	{	\left\{
			\begin{array}{ll}
				#1 & \mbox{if } #2 \\
				#3 & \mbox{if } #4
			\end{array}
		\right.}
%	\newcommand{\threepartdef}[6]
%	{	\left\{
%			\begin{array}{ll}
%				#1 & \mbox{if } #2 \\
%				#3 & \mbox{if } #4 \\
%				#5 & \mbox{if } #6
%			\end{array}
%		\right.}
%Blackboard bold	
	\newcommand{\N}{{\mathbb N}} 
	\newcommand{\R}{{\mathbb R}} 
	\newcommand{\C}{{\mathbb C}} 
	%Curly letters
%	\newcommand{\E}{{\mathcal E}}	
%	\newcommand{\F}{{\mathcal F}} 
%	\newcommand{\Null}{{\mathcal N}} 
%	\newcommand{\R}{\mathcal{R}} 
%	\newcommand{\Z}{\mathcal{Z}} 
%Uprights
	\newcommand{\D}{{\mathrm D}}
	\renewcommand{\d}{{\mathrm d}}
%Easy Greeks
%	\def\a{\alpha}
%	\def\g{\gamma}
%	\newcommand{\e}{{\varepsilon}}
%	\newcommand{\la}{{\lambda}}  
%	\newcommand{\om}{{\Omega}}
	\renewcommand{\epsilon}{\varepsilon}
	
%Vector operators
	\def \p {\partial}
	\newcommand{\del}{{\nabla}}
	\newcommand{\grad}{{\nabla}}
	\newcommand{\cross}{{\times}}
	\newcommand{\Lap}{{\nabla^2}} 
	\newcommand{\vnabla}{\boldsymbol{\nabla}}
%Arrows
%	\newcommand{\goesto}[1]{\xrightarrow[#1]{}}
%hats
%	\newcommand{\nhat}{\widehat{\v{n}}}
%	\newcommand{\shat}{\widehat{\v{s}}}
	\renewcommand{\hat}{\widehat}
%Note to self
%	\newcommand{\note}[1]{[\textbf{\textsl{\MakeUppercase{#1}}}]}
	
%Stress tensor
%	\newcommand{\T}{\vg{\sigma}}

% Vectors, quaternions etc.
\renewcommand{\v}[1]{\mathbfit{#1}}      % Generic vector
\newcommand{\vc}[1]{\bm{\mathcal{#1}}}      % vector mathcal
\newcommand{\uv}[1]{\widehat{\mathbfit{#1}}}      % Generic unit vector
\newcommand{\q}[1]{\mathsfbfit{#1}}        % Generic quaternion
\renewcommand{\t}[1]{\mathsfbfit{#1}}	 % Generic tensor
\newcommand{\m}[1]{\mathsfbfit{#1}}	 % Generic matrix
\newcommand{\vu}[1]{\mathbf{#1}}         % Upright vector (for zero vector)
\newcommand{\qu}[1]{\mathbf{#1}}       % Upright quaternion (for zero quaternion)
\newcommand{\tu}[1]{\bm{\mathsf{#1}}}	 % Upright tensor (for zero tensor)

%Real and imaginary
\renewcommand{\Re}{\operatorname{Re}}
\renewcommand{\Im}{\operatorname{Im}}
	
%Non-dim numbers
%	\renewcommand{\Re}{\operatorname{Re}}
%	\newcommand{\Bi}{\operatorname{Bi}}
%	\newcommand{\Fo}{\operatorname{Fo}}
	
%Misc
\DeclareMathSymbol{\Delta}{\mathalpha}{operators}{1} % Keeps Delta upright
\DeclareMathSymbol{\Gamma}{\mathalpha}{operators}{0} % Keeps Gamma upright
	\newcommand{\cosec}{\operatorname{cosec}}
	\newcommand{\cosech}{\operatorname{cosech}}
	\newcommand{\sech}{\operatorname{sech}}
	\newcommand{\tr}{\operatorname{tr}}
	\newcommand{\erf}{\operatorname{erf}}
	\newcommand{\order}{\mathcal{O}}
%	\newcommand{\V}{\mathcal{V}} % 	CurlyV = (Bi V) / (1 + Bi)

\newcommand{\mvec}[2]{\begin{pmatrix}#1\\#2\end{pmatrix}}
\newcommand{\mat}[4]{\begin{pmatrix}#1&#2\\#3&#4\end{pmatrix}}


% Differentiating 
\newcommand{\fd}[2]{\mathchoice{\frac{\d #1}{\d #2}}{\d #1/\d #2}{\d #1/\d #2}{\d #1/\d #2}}
\newcommand{\fdd}[2]{\mathchoice{\frac{\d^2 #1}{\d #2^2}}{\d^2 #1/\d #2^2}{\d^2 #1/\d #2^2}{\d^2 #1/\d #2^2}}
\newcommand{\fddd}[2]{\mathchoice{\frac{\d^3 #1}{\d #2^3}}{\d^3 #1/\d #2^3}{\d^3 #1/\d #2^3}{\d^3 #1/\d #2^3}}
\newcommand{\fdddd}[2]{\mathchoice{\frac{\d^4 #1}{\d #2^4}}{\d^4 #1/\d #2^4}{\d^4 #1/\d #2^4}{\d^4 #1/\d #2^4}}
\newcommand{\fdn}[3]{\mathchoice{\frac{\d^{#1} #2}{\d #3^{#1}}}{\d^{#1} #2/\d #3^{#1}}{\d^{#1} #2/\d #3^{#1}}{\d^{#1} #2/\d #3^{#1}}}
\newcommand{\pd}[2]{\mathchoice{\frac{\p #1}{\p #2}}{\p #1/\p #2}{\p #1/\p #2}{\p #1/\p #2}}
\newcommand{\pdd}[2]{\mathchoice{\frac{\p^2 #1}{\p #2^2}}{\p^2 #1/\p #2^2}{\p^2 #1/\p #2^2}{\p^2 #1/\p #2^2}}
\newcommand{\pddmixed}[3]{\frac{\p^2 #1}{\p #2 \p #3}}
\newcommand{\pddd}[2]{\frac{\p^3 #1}{\p #2^3}}
\newcommand{\pdn}[3]{\mathchoice{\frac{\p^{#1} #2}{\p #3^{#1}}}{\p^{#1} #2/\p #3^{#1}}{\p^{#1} #2/\p #3^{#1}}{\p^{#1} #2/\p #3^{#1}}}
	
	% Misc
\newcommand{\e}{\mathrm{e}}
\renewcommand{\i}{\mathrm{i}}
\newcommand{\dx}{\,\mathrm{d}x}
\newcommand{\dy}{\,\mathrm{d}y}
\renewcommand{\epsilon}{\varepsilon}
\renewcommand{\Re}{\operatorname{Re}}
\newcommand{\Four}{\mathcal{F}}
	\definecolor{nem-red}{HTML}{FFEDF0}
\definecolor{nem-dark-red}{HTML}{d21f3c}
	\newenvironment{nem}{\begin{tcolorbox}[colback=nem-red,colframe=nem-dark-red,breakable,enhanced,parbox=false,title={\sf\small{Non-examinable}}]}{\end{tcolorbox}}
	
\newenvironment{colour-box-breakable}[3]{
\setlist[itemize]{topsep=-\parskip,before=\vspace{-2mm},itemsep=3pt}
%\setlength{\parskip}{0pt}
\begin{tcolorbox}[
	enhanced,
    breakable,
	colback=#2,
	colframe=#1,
	colbacktitle=#2, 
	frame hidden,
    boxrule=0pt,
    leftrule=2pt,
	toptitle=2.5mm,
	bottomtitle=-.5mm, 
    pad before break=1pt,
	pad after break=6pt,
	parbox=false,
	bottom=10pt,
	%grow to right by=0.5mm,
    %grow to left by=-0.2mm,
	sharp corners=west,
  	fonttitle={\sf\small},
	fontupper={\small},
	title={\color{#1}#3}, 
  	overlay first={
	  \fill[#1]
      ([xshift=2pt]frame.north west) --
      ([xshift=2pt]frame.south west) --
      ([xshift=0pt]frame.south west) --
      ([yshift=-2pt]frame.north west) arc
      [radius=2pt,start angle=180,end angle=90]--
      cycle;
    },
     overlay unbroken={
	  \fill[#1]
      ([xshift=2pt]frame.north west) --
      ([xshift=2pt]frame.south west) arc
      [radius=2pt,start angle=-90,end angle=-180] --
      ([yshift=-2pt]frame.north west) arc
      [radius=2pt,start angle=180,end angle=90]--
      cycle;
    },
  	overlay last={
	  \fill[#1]
      ([xshift=2pt]frame.north west) --
      ([xshift=2pt]frame.south west) arc
      [radius=2pt,start angle=-90,end angle=-180] --
      ([yshift=-2pt]frame.north west) --
      ([yshift=0pt]frame.north west) --
      cycle;
    },
]}{\end{tcolorbox}}

\definecolor{commentary-green}{HTML}{FBF3EE}
\definecolor{commentary-dark-green}{HTML}{c57a44}

\newcommand{\commentary}[1]
{
\ifbool{show_commentary}{
	\begin{colour-box-breakable}
		{commentary-dark-green}
		{commentary-green}
		{\footnotesize\raisebox{-2pt}{\includegraphics[width=1em]{../common/squirrel.png}}\,\,Feedback}
	#1
\end{colour-box-breakable}
	}{}
}
%}{
%  \newenvironment{commentary}{\commentaryA}{\endcommentary}
%}

%Depth of display breaks to allow	
	\allowdisplaybreaks[1]

%How deep do we want the Table of Contents to go?	
%	\setcounter{tocdepth}{1}

%Sections
\makeatletter
\renewcommand{\section}{\@startsection
{section}%                   % the name
{1}%                         % the level
{0mm}%                       % the indent
{-\baselineskip}%            % the before skip
{0.5\baselineskip}%          % the after skip
{\normalfont\bfseries}} % the style
\makeatother

% Wavy underline
\makeatletter
\font\sixly=lasyb10 scaled 1200
\def\tinners{\bgroup \markoverwith{\lower8.5\p@\hbox{\sixly \char58}}\ULon}
\makeatother
\newcommand{\answer}[1]{\tinners{\; #1 \;}}

%Setting the footnotes to use *,**,+ etc rather than 1,2,3	
\renewcommand{\thefootnote}{\fnsymbol{footnote}}

%Italic quote environment
%\newenvironment{quoteitalic}{\begin{quote}\itshape}{\end{quote}}

%Heading
\newcommand{\heading}[2]{
\setlength{\parindent}{0pt}
\textbf{MATH3171 (Mathematical Biology III)}

\textbf{YEAR 2026--27, \MakeUppercase{#1} TERM}
\setlength{\parskip}{3ex}

\textbf{\MakeUppercase{#2}}
%
\setlength{\parskip}{2ex}

}

%========================================================================================================================
%DOCUMENT BEGINS HERE!

\usepackage{booktabs}

\setbool{show_question}{true}
\setbool{show_solution}{false}
\setbool{show_commentary}{false}

\begin{document}

\heading{Epiphany}{Problem sheet 7\solutioninheading}

\renewcommand*{\thefootnote}{\arabic{footnote}}

\ifbool{show_solution}{}{
    This problem sheet is due in by \textbf{noon} on
    \textbf{Monday X X 2026}.

    \textbf{Topics covered}: Chapters 12--13
}

\begin{enumerate}

    %------------------------------------------------------------------%
    % Q1
    %------------------------------------------------------------------%

    \question{}{
        \textbf{Chemotaxis in a Petri dish}

        A chemotactic bacterium is kept in a circular Petri dish, and
        moves according to self-generated gradients of a chemoattractant.
        A model for this on the timescale neglecting bacterial growth or
        death is given by
        \begin{align*}
            \pd{n}{t}
            &= D\nabla^2 n
            - \vnabla\cdot\bigl(g(n,c)\vnabla c\bigr), \\
            \pd{c}{t}
            &= \nabla^2 c+f(n)-ac,
        \end{align*}
        where $a,D$ are positive parameters, and the solution is solved
        in the domain (using polar coordinates $(r,\theta)$)
        $0\leq r\leq R$ and $0\leq\theta<2\pi$. We assume both species
        satisfy no-flux conditions on the domain boundary $r=R$.

        \begin{enumerate}[label=(\alph*)]

            % (a)
            \item Give a biological interpretation to the parameters as
            well as the two arbitrary functions.

            % (b)
            \item Is the no-flux boundary condition on both species the
            same as Neumann boundary conditions? Clearly show why or why
            not.

            % (c)
            \item Find all spatially homogeneous equilibria of this model,
            and state any conditions necessary for these equilibria to be
            stable in the absence of transport terms. NB: perturbation
            growth rates with $\lambda=0$ in this model can be considered
            as stable, as corresponding perturbations will change the
            total mass of the system.

            % (d)
            \item Linearise the PDE system, and determine an equation for
            the growth rates of a perturbation, $\lambda_k$, in terms of
            spatial eigenvalues $\rho_k$ of the Laplacian (you do not need
            to solve for $\lambda_k$ but should have an explicit algebraic
            equation that it satisfies).

            % (e)
            \item \emph{Optional.} What equation and boundary conditions
            determine the spatial eigenvalues $\rho_k$? You do not need
            to solve this spatial eigenvalue problem, but should describe
            (in words or pictures) what kinds of functions satisfy it.

        \end{enumerate}
    }{
        \begin{enumerate}[label=(\alph*)]

            % (a)
            \item
            $D$ is a ratio of diffusion coefficients (cell population to
            chemoattractant), $a$ is a decay rate of the chemoattractant,
            $g$ is the response of the cell population to the
            chemoattractant, and $f$ is the production of $c$ by the cells.

            \commentary{
                Broadly done well. Note that $D$, $f$, $g$, $h$ all have
                some units of $1/\text{time}$, which are `rates.'
                Similarly, note that the diffusion constant $D$ is
                partially non-dimensionalised in that the space variable
                has been scaled such that $c$ has a diffusion rate of 1.
                This isn't crucial, but is a slightly more nuanced
                understanding of how these terms can be interpreted.
            }

            %------------------------------------------------------------------%

            % (b)
            \item
            Yes. We have
            \[
                \underbrace{
                    D\pd{n}{r}-g(n,c)\pd{c}{r}=0,
                    \qquad
                    \pd{c}{r}=0
                }_{\textrm{No-flux}}
                \iff
                \underbrace{
                    D\pd{n}{r}=0,
                    \qquad
                    \pd{c}{r}=0
                }_{\textrm{Neumann}},
            \]
            as the equivalence we want to show, with all quantities
            evaluated at the boundary $r=R$. The $\impliedby$ direction
            is direct, whereas the $\implies$ direction just requires
            noticing that the no-flux term on $c$ forces the no-flux
            term on $n$ to imply Neumann boundary conditions on $n$.

            \commentary{
                Quite a few conceptual errors here. It is important to
                just explain in words and a few symbols why one set of
                boundary conditions implies the other, and vice-versa.
                In general, no-flux conditions will not be the same, but
                in this case they are the same due to no-flux on $c$
                simplifying the flux of $n$.
            }

            %------------------------------------------------------------------%

            % (c)
            \item
            Homogeneous steady states satisfy
            \[
                f(n_0)-ac_0=0,
            \]
            but we do not have a second equation, and so this is true
            for any value of $n_0\geq0$.

            Linearising about such a steady state, we find the Jacobian
            \[
                \m{J}
                =
                \begin{pmatrix}
                    0 & 0 \\
                    f'(n_0) & -a
                \end{pmatrix},
            \]
            and hence the growth rates of perturbations are
            $\lambda=0,-a$. Hence, any such equilibrium is stable except
            for perturbations that change the total mass of cells.

            \commentary{
                Generally fine solutions. It can help to explain that
                there are infinitely many equilibria given by some
                $n_0\geq0$, and they are all stable in the sense that if
                you perturb only $c_0$, those perturbations will decay.
                Perturbations in $n_0$ will change its value, but they
                won't grow or decay.
            }

            %------------------------------------------------------------------%

            % (d)
            \item
            Perturbations $(n_1,c_1)$ satisfy the linear system
            \begin{align*}
                \pd{n_1}{t}
                &=D\nabla^2n_1-g(n_0,c_0)\nabla^2c_1, \\
                \pd{c_1}{t}
                &=\nabla^2c_1+f'(n_0)n_1-ac_1.
            \end{align*}
            Substituting in
            \[
                (n_1,c_1)\propto
                \e^{\lambda_k t}w_k(r,\theta),
            \]
            and dividing out this factor, we find that the $\lambda_k$
            are eigenvalues of the matrix
            \[
                \m{M}
                =
                \begin{pmatrix}
                    -D\rho_k & g(n_0,c_0)\rho_k \\
                    f'(n_0) & -\rho_k-a
                \end{pmatrix}.
            \]
            So that $\lambda_k$ satisfies the quadratic
            \[
                \lambda_k^2
                +\bigl((D+1)\rho_k+a\bigr)\lambda_k
                +\bigl(D\rho_k^2+aD\rho_k
                -g(n_0,c_0)\rho_k f'(n_0)\bigr)
                =0.
            \]
            Alternatively, these terms can be written as the trace and
            determinant of $\m{M}$.

            \commentary{
                This was done rather well by most people who attempted it.
            }

            %------------------------------------------------------------------%

            % (e)
            \item
            The spatial eigenvalue problem reads
            \[
                \nabla^2w_k+\rho_kw_k=0,
            \]
            which has solutions in terms of $\sin$ and $\cos$ in the
            $\theta$ direction, and Bessel functions in the $r$
            direction. The radial component has to satisfy Neumann
            boundary conditions, which are used to determine the value
            of $\rho_k$.

        \end{enumerate}
    }

    %------------------------------------------------------------------%
    % Q2
    %------------------------------------------------------------------%

    \question{}{
        \textbf{Colony formation in an ambush predator--prey system}

        Consider the following reaction--advection--diffusion system,
        \begin{equation}
            \begin{aligned}
                \label{ambush_model}
                \pd{u}{t}
                &=u(1-u)-uv+\nabla^2u
                -\alpha\vnabla\cdot\left(u\vnabla v\right), \\
                \pd{v}{t}
                &=-cv+uv+D\nabla^2v,
            \end{aligned}
        \end{equation}
        where all parameters are positive, and we have chosen a timescale
        such that $D$ represents the ratio of diffusion parameters. The
        term involving $\alpha$ represents movement of the prey species
        $u$ \emph{towards} the predator, so that it is attracted to the
        predator. This occurs in nature for many examples of
        \emph{ambush} predators, as noted in the lecture notes.

        \begin{enumerate}[label=(\alph*)]

            % (a)
            \item Compute the homogeneous equilibrium where both species
            coexist, and determine any conditions necessary for its
            feasibility and stability. Can this Jacobian satisfy the
            Turing conditions if we only had reaction and diffusion
            (that is, if $\alpha=0$)?

            % (b)
            \item In the lecture notes we derived the following two
            necessary conditions for diffusion-driven (Turing)
            instability in cross-diffusion systems:
            \[
                d_1G_v+d_4F_u>d_2G_u+d_3F_v,
            \]
            \[
                (d_1G_v+d_4F_u-d_2G_u-d_3F_v)^2
                >4(d_1d_4-d_2d_3)(F_uG_v-F_vG_u).
            \]
            Write the functions $d_i$ for the system
            \eqref{ambush_model}, and hence find conditions for this
            system to admit Turing instabilities. Simplify these
            conditions to find that the parameters need to satisfy
            \[
                \alpha>
                \frac{D}{1-c}
                +2\sqrt{\frac{D}{c(1-c)}}
            \]
            and briefly interpret how $\alpha$ and $D$ influence the
            system's ability to form patterns.

        \end{enumerate}
    }{
        \begin{enumerate}[label=(\alph*)]

            % (a)
            \item
            For both species to exist we need $u_0\neq0$ and
            $v_0\neq0$, so we have that an equilibrium must satisfy
            \[
                1-u_0-v_0=0,
                \qquad
                -c+u_0=0.
            \]
            Hence, $u_0=c$ and $v_0=1-c$ is our coexistence state.

            The Jacobian at this steady state is given by
            \[
                \m{J}
                =
                \begin{pmatrix}
                    1-2u_0-v_0 & -u_0 \\
                    v_0 & -c+u_0
                \end{pmatrix}
                =
                \begin{pmatrix}
                    -c & -c \\
                    1-c & 0
                \end{pmatrix}.
            \]
            To have $v_0>0$ we need $c<1$. This is also exactly the
            condition we need for stability. One way to quickly see this
            is to note that
            \[
                \tr(\m{J})=-c<0
            \]
            (which is always satisfied) and
            \[
                \det(\m{J})=c(1-c)>0
            \]
            if $c<1$.

            As the Jacobian has a zero, it cannot satisfy the Turing
            conditions. Hence, if there is any pattern formation in the
            system, we will require $\alpha\neq0$.

            \commentary{
                Remember that in the purely reaction--diffusion case,
                you have constraints on the Jacobian that can ever
                satisfy the Turing conditions (Section 11.3). These can
                be quoted directly.
            }

            %------------------------------------------------------------------%

            % (b)
            \item
            We have already computed the Jacobian evaluated at the
            steady state. If we linearise the system, using the notation
            from the lectures, we find
            \[
                d_1=1,
                \qquad
                d_2=-\alpha u_0=-\alpha c,
                \qquad
                d_3=0,
                \qquad
                d_4=D.
            \]
            This leads to the following two conditions for instability:
            \[
                -Dc>-\alpha c(1-c),
                \qquad
                \bigl(-Dc+\alpha c(1-c)\bigr)^2
                >4Dc(1-c).
            \]
            Simplifying these, we have
            \[
                D<\alpha(1-c),
                \qquad
                \bigl(\alpha(1-c)-D\bigr)^2
                >\frac{4D(1-c)}{c}.
            \]
            We can divide the first inequality by $(1-c)$ and the
            second by $(1-c)^2$ to get
            \[
                \frac{D}{1-c}<\alpha,
                \qquad
                \left(\alpha-\frac{D}{1-c}\right)^2
                >\frac{4D}{c(1-c)}.
            \]
            The first inequality can be thought of as requiring us to
            take the positive branch of the square root in the second.
            Doing this, we can solve for $\alpha$ as
            \[
                \alpha>
                \frac{D}{1-c}
                +2\sqrt{\frac{D}{c(1-c)}}.
            \]

            We see that, for fixed $D$ and $c$, $\alpha$ can always be
            increased to lead to a spatial instability. In contrast,
            for fixed $\alpha$ and $c$, $D$ cannot be increased above a
            threshold before the system loses the ability to pattern,
            though decreasing it can always allow for pattern formation
            if $\alpha$ is nonzero.

            \commentary{
                Generally fine. It can help to explain any algebra you
                do, e.g. why you take the positive branch when solving
                the quadratic in $D$. Clear written explanations are
                very valuable for understanding your thinking.
            }

        \end{enumerate}
    }

\end{enumerate}

\end{document}
